% Part: many-valued-logic
% Chapter: syntax-and-semantics
% Section: valuations-sat

\documentclass[../../../include/open-logic-section]{subfiles}

\begin{document}

\olfileid{mvl}{syn}{val}

\olsection{\usetoken{P}{valuation} અને સંતોષ}

\begin{defn}[!!^{valuation}s]
$V$ સત્યમૂલ્યોનો ગણ હોય. $\Lang{L}$ માટે $V$માં મૂલ્યો ધરાવતું
\emph{!!{valuation}} એ ભાષાના !!{propositional variable}sને $V$નો
!!a{element} નિયુક્ત કરતું વિધેય~$\pAssign{v}$ છે; એટલે કે,
$\pAssign{v} \colon \PVar \to V$.
\end{defn}

\begin{defn}\ollabel{defn:pValue}
બહુમૂલ્ય તર્કશાસ્ત્ર~$\Log L$ના સત્યમૂલ્યોના ગણ~$V$માં મૂલ્યો
ધરાવતું !!a{valuation}~$\pAssign{v}$ આપેલું હોય, ત્યારે મૂલ્યાંકન
વિધેય $\pValue{v} \colon \Frm[L] \to V$ને નીચે પ્રમાણે અનુમાનાત્મક
રીતે વ્યાખ્યાયિત કરો:
  \begin{enumerate}
    \item $\pValue{v}(\Obj p_n) = \pAssign{v}(\Obj p_n)$;
    \item જો $\star$ એ $0$-સ્થાનીય સંયોજક હોય, તો $\pValue{v}(\star)  = \tf{\star}[\Log L]$;
    \item જો $\star$ એ $n$-સ્થાનીય સંયોજક હોય, તો
    \[
      \pValue{v}(\star(!A_1, \dots, !A_n)) = \tf{\star}[\Log L]
      (\pValue{v}(!A_1), \dots, \pValue{v}(!A_n)).
    \]
  \end{enumerate}
\end{defn}

\begin{defn}[સંતોષ]
\ollabel{defn:satisfaction} !!{formula}~$!A$નો !!a{valuation}~$\pAssign{v}$
વડે \emph{સંતોષ} થાય છે, $\pSat{v}{!A}[\Log L]$, ત્યારે જ જ્યારે
$\pValue{v}(!A)[\Log L] \in V^+$ હોય; અહીં $V^+$ એ~$\Log L$ના
નિર્દિષ્ટ સત્યમૂલ્યોનો ગણ છે.

``$\pSat{v}{!A}[\Log L]$ નથી'' એવો અર્થ દર્શાવવા આપણે
$\pSat/{v}{!A}[\Log L]$ લખીએ છીએ. જો $\Gamma$ !!{formula}sનો ગણ હોય, તો
દરેક~$!A \in \Gamma$ માટે $\pSat{v}{!A}[\Log L]$ હોય ત્યારે જ
$\pSat{v}{\Gamma}[\Log L]$.
\end{defn}

\end{document}
